% Result-dependent sentences. Every number below is copied from the generated tables
% (tables/main_*_k25.tex, review_cost.tex, significance.tex, recog_*.tex, mp_analysis.tex,
% font_sensitivity.tex, char_error_subset.tex, audit.tex) after all runs finished.
\newcommand{\SROIESENTENCE}{On held-out SROIE the ranking of baselines flips: TrOCR's printed checkpoints were
fine-tuned on SROIE receipts, so TrOCR re-recognition is strong (.979 macro-F1) while PARSeq agreement fails on these
long receipt lines (.742); \ours{} is slightly below the former in macro-F1 (.973) and above it in AUROC (.992 vs.\ .978,
Table~\ref{tab:significance}). No single re-recognition baseline is the best choice on every dataset (PARSeq on forms
and scene text, TrOCR on SROIE, neither on handwriting), whereas \ours{} is within 0.006 macro-F1 of the best baseline
or better on four of the five datasets. Its average macro-F1 (.954) is the highest of all methods in
Table~\ref{tab:main}: well above the best single baseline (.907, PARSeq), equal to choosing the best baseline separately
for every dataset in hindsight (.954), and marginally above simple learners trained on our own counterfactual evidence
(.939--.952).}
\newcommand{\REVIEWSENTENCE}{At 95\% accepted precision, \ours{} requires reviewing 44.7\% of SROIE lines (TrOCR
re-recognition 45.0\%, PARSeq 69.5\%, the VLM judge 76.2\%, engine confidence 90.9\%) and 51.4\% of IAM words (PARSeq
54.6\%, the VLM judge 94.8\%). On FUNSD, CORD and IIIT5K, \ours{} and PARSeq are within two points of each other (FUNSD
37.2\% vs.\ 35.4\%, CORD 37.4\% vs.\ 38.6\%, IIIT5K 25.6\% vs.\ 25.3\%), while engine confidence and the VLM judge require
1.6--2.3$\times$ more review than \ours{} on the three receipt and form datasets. At the strictest target (98\%) PARSeq
is better on FUNSD and IAM.}
\newcommand{\XDSSENTENCE}{The ranking transfers well across all twenty dataset pairs (AUROC .892--.995;
e.g.\ FUNSD$\rightarrow$CORD .995 vs.\ .996 in-domain, CORD$\rightarrow$FUNSD .968 vs.\ .979). The weakest transfers are
those \emph{to} handwriting (SROIE$\rightarrow$IAM .892, CORD$\rightarrow$IAM .954) and from receipt lines to forms
(SROIE$\rightarrow$FUNSD .941). As for engines, macro-F1 at the fixed 0.5 threshold transfers much less well than the
ranking (e.g.\ SROIE$\rightarrow$FUNSD .620).}
\newcommand{\SIGSENTENCE}{Against the strongest baseline, \ours{} is significantly better on IAM in both metrics
($\Delta$AUROC $+0.055$, $\Delta$F1 $+0.026$, $p<0.001$), better in AUROC on CORD at the document level ($+0.007$,
$p=0.023$; macro-F1 $+0.005$, not significant) and on SROIE ($+0.014$, $p<0.001$, while TrOCR re-recognition is
slightly better in macro-F1, $-0.006$, $p=0.002$); it is statistically tied with PARSeq on FUNSD at the document level
($p=0.25$ and $p=0.12$) and significantly worse than PARSeq on IIIT5K ($-0.004$ AUROC, $-0.021$ macro-F1). It is
significantly better than the VLM judge and than TrOCR re-recognition on every dataset except SROIE, where the latter is
in-domain.}
\newcommand{\MPSENTENCE}{Three observations explain the null effect of message passing. First, removing message
passing changes neither metric by more than 0.009 on any dataset (A vs.\ B). Second, the \emph{explicit} comparison of each
rival with the claim does matter: removing the claim-relative features costs 0.010 to 0.023 macro-F1 on CORD, SROIE and
IIIT5K (A vs.\ C), and adding message passing back does not recover it (C vs.\ D). The network is not learning the
comparison from the graph; it uses the comparison when it is computed for it. Third, the graph is a depth-one star whose
readout already pools over every rival, so a message-passing layer can only re-express information the readout sees
anyway. Removing the rivals entirely (E) costs ranking quality on FUNSD (AUROC .974 to .961) but little elsewhere,
because the claim node keeps its own evidence and its agreement with both verifiers' readings, which summarise what the
strongest rivals would say. In short, what helps is counterfactual \emph{evidence}: scoring the claim against
alternatives under prior-cancelled likelihoods and making the comparison explicit; the graph is a convenient container
for it, not the source of the gain.}
\newcommand{\FONTSENTENCE}{The trained model is insensitive to the rendering: across seven
renderings its macro-F1 varies by at most 0.005 on every dataset (FUNSD .937--.939, CORD .972--.977, IAM .919--.923), and
the AUROC of the claim's ratio alone by at most 0.011. The training-free rule is more sensitive (AUROC .920--.940 on FUNSD,
.919--.944 on IAM); it is best with a serif or monospaced font rather than our default sans-serif, which we did not tune, and
a handwriting-like font does not help on handwriting. A mismatched rendering is thus absorbed by learning; only the
label-free rule would benefit from choosing the font.}
\newcommand{\CHARSENTENCE}{The removed share is large on FUNSD (51\% of errors) and moderate on CORD, SROIE and IAM
(18--23\%); IIIT5K's canonical protocol already excludes such differences. \ours{}' AUROC changes by at most 0.011
(it rises slightly or is unchanged), and its ordering relative to the baselines is unchanged except that the VLM judge,
whose AUROC rises most (FUNSD .846 to .928, SROIE .941 to .982), is visibly penalised by these artefacts. The conclusions of
Sec.~\ref{sec:main_results} therefore do not rest on annotation conventions.}
\newcommand{\AUDITSENTENCE}{Two conclusions follow. First, fewer than half of \ours{}' apparent
mistakes are genuine verification errors on a correct label (42 of 100 sampled): 25 are label noise, 9 crop problems and
8 ambiguous, and on IIIT5K 16 of 20 are claims that count as correct only after the case-insensitive, punctuation-free
normalisation of the protocol (\emph{e.g.}\ \texttt{HoLLywooD} for an upper-case sign), which \ours{}, judging the exact
appearance of the string, rejects; this protocol mismatch, rather than a visual error, explains most of \ours{}' gap to
PARSeq agreement (which compares normalised strings) on IIIT5K. Second, the stratum ``errors that PARSeq agrees with''
consists mostly of cases where the claim matches what is visible: only 10 of 92 sampled words are genuine errors, against
46 label-noise and 21 crop problems, so low rejection rates there are not a weakness of either verifier. The genuine
failures that remain fall into two identifiable modes: rejecting correct short or case-ambiguous claims when PARSeq reads
another case (CORD: \texttt{Rp}, \texttt{x}, \texttt{Pop}), and accepting capitalisation errors in handwriting (IAM:
\texttt{pinkie}, \texttt{market}, \texttt{church}).}
